% ====================================================================
% 第3章：方法分类（LaTeX版本，含TikZ图）
% ====================================================================

\section{方法分类}\label{ch:methods}

\subsection{本章概述}

Agent Evolution 方法众多，名称也不统一。本章按"主要选择压力与主要可变对象"进行分类：基于奖励的进化主要优化策略或行为分布；自博弈进化主要生成任务、对手和验证信号；提示词进化主要改变上下文、反思、原则和输出修订流程；架构搜索主要改变代码、工具、控制流和多智能体拓扑；基于记忆的进化主要改变长期经验、技能库和检索策略；混合方法则把多种机制组合成端到端闭环。

\begin{figure}[htbp]
\centering
\begin{tikzpicture}[
  level 1/.style={sibling distance=42mm, level distance=22mm},
  level 2/.style={sibling distance=22mm, level distance=20mm},
  root/.style={rectangle, rounded corners, draw=cat1, fill=cat1!10,
    text width=60mm, align=center, font=\bfseries, minimum height=10mm},
  cat/.style={rectangle, rounded corners, draw=#1, fill=#1!8,
    text width=30mm, align=center, font=\small\bfseries, minimum height=8mm},
  paper/.style={rectangle, rounded corners, draw=gray!60, fill=gray!5,
    text width=26mm, align=left, font=\scriptsize, inner sep=2pt},
  edge from parent/.style={draw, -Stealth, thick, gray!70}
]
\node[root] {Agent Self-Evolution\\方法分类}
  child {
    node[cat=cat5] {3.1 基于奖励\\的进化}
    child { node[paper] {STaR, Self-Rewarding\\Meta-Rewarding\\RISE, Agent-R} }
    child { node[paper] {RAGEN, IterAlign\\SICA, DGM} }
  }
  child {
    node[cat=cat2] {3.2 自博弈\\进化}
    child { node[paper] {Absolute Zero\\SPIRAL} }
    child { node[paper] {Multi-Agent Debate\\Self-Play FT} }
  }
  child {
    node[cat=cat3] {3.3 提示词\\进化}
    child { node[paper] {Reflexion, ExpeL\\Self-Refine} }
    child { node[paper] {ACE, ICE, EvolveR\\ReasoningBank} }
  }
  child {
    node[cat=cat4] {3.4 架构\\搜索}
    child { node[paper] {ADAS, Gödel Agent\\DGM, SICA} }
    child { node[paper] {EvoMAC, AlphaEvolve\\FunSearch} }
  }
  child {
    node[cat=cat5] {3.5 记忆\\进化}
    child { node[paper] {Voyager, Memory-R1\\AriadneMem} }
    child { node[paper] {SEDM, A-Mem\\Memento} }
  }
  child {
    node[cat=cat6] {3.6 混合\\方法}
    child { node[paper] {EvolveR (全生命周期)\\WebEvolver} }
    child { node[paper] {HyperAgents\\EvoAgentX} }
  };
\end{tikzpicture}
\caption{Agent Self-Evolution 方法分类体系。六大类别按主要选择压力和可变对象划分。}
\label{fig:ch3-taxonomy}
\end{figure}

\subsection{基于奖励的进化}\label{sec:reward}

基于奖励的进化是最接近传统强化学习和后训练的一类方法。核心思想是：把智能体交互轨迹、推理过程或输出结果映射为标量或偏好信号，再通过策略优化、偏好学习、筛选微调或archive selection使后续行为更可能获得高奖励。

\subsubsection{代表系统}

\textbf{STaR}（Self-Taught Reasoner）是早期代表。它把"模型自己生成可学习材料"组织成闭环：生成、验证、训练、再生成。虽然它主要作用于推理链和模型权重，但其bootstrapping逻辑成为后续自演化研究的基础模板。

\textbf{Self-Rewarding Language Models} 把奖励信号进一步内生化。系统让同一个模型既生成候选回答，又作为judge评价候选，并用Iterative DPO进行多轮偏好优化。\textbf{Meta-Rewarding}进一步区分actor、judge和meta-judge，让模型不仅评价回答，还评价评价本身。

\textbf{RISE}将单轮问题改写为多轮MDP，让模型学习在测试时递归自省和修正。其关键是把"再想一轮"从prompt trick变成可训练策略。

\begin{figure}[htbp]
\centering
\begin{tikzpicture}[
  box/.style={rectangle, rounded corners=3pt, draw=#1, fill=#1!10,
    text width=35mm, align=center, font=\small, minimum height=12mm},
  arr/.style={-Stealth, thick, gray!60}
]
\node[box=cat1] (gen) at (0,0) {生成候选\\(Actor)};
\node[box=cat2] (judge) at (4.5,0) {自评/外部评估\\(Judge/Meta-Judge)};
\node[box=cat3] (train) at (9,0) {偏好优化\\(DPO/GRPO)};
\node[box=cat4] (iter) at (4.5,-2.2) {迭代更新\\(权重/策略)};

\draw[arr] (gen) -- (judge);
\draw[arr] (judge) -- (train);
\draw[arr] (train) -- (iter);
\draw[arr] (iter.west) -| (gen.south);

\node[font=\scriptsize, text=gray] at (2.2, 0.6) {候选+分数};
\node[font=\scriptsize, text=gray] at (6.7, 0.6) {偏好对};
\node[font=\scriptsize, text=gray] at (4.5, -1.0) {更新后模型};
\end{tikzpicture}
\caption{基于奖励的进化通用闭环：Actor生成候选，Judge（自评或外部）评分，DPO/GRPO优化偏好，迭代更新模型权重。}
\label{fig:ch3-reward-loop}
\end{figure}

\subsubsection{核心挑战}

基于奖励的方法面临\textbf{Goodhart风险}：模型可能学会迎合评分规则而非真正提升能力。\textbf{样本效率}问题突出——RL训练需要大量交互轨迹，成本高昂。\textbf{泛化不确定}——在特定benchmark上的改进不一定迁移到真实场景。

\subsection{自博弈进化}\label{sec:selfplay}

自博弈进化把"数据从哪里来"这个问题推到系统内部。\textbf{Absolute Zero}提出在零外部数据条件下，由同一个模型提出最大化自身学习进度的任务并尝试解决。\textbf{SPIRAL}让语言模型在零和游戏中自我博弈，发现可迁移的推理能力。

\begin{figure}[htbp]
\centering
\begin{tikzpicture}[
  role/.style={circle, draw=#1, fill=#1!15, minimum size=16mm,
    font=\small\bfseries, align=center},
  arr/.style={-Stealth, thick, #1}
]
\node[role=cat2] (prop) at (0,0) {Proposer\\(出题)};
\node[role=cat1] (solve) at (5,0) {Solver\\(解题)};
\node[role=cat3] (verif) at (2.5,-2.5) {Verifier\\(验证)};

\draw[arr=cat2] (prop) -- node[above, font=\scriptsize]{任务} (solve);
\draw[arr=cat1] (solve) -- node[right, font=\scriptsize, align=left]{解答\\+反馈} (verif);
\draw[arr=cat3] (verif) -- node[left, font=\scriptsize, align=right]{学习\\进度信号} (prop);

\node[font=\scriptsize, text=gray, align=center] at (2.5, 1.2)
  {关键：任务难度需位于能力边界\\太易无训练价值，太难无有效反馈};
\end{tikzpicture}
\caption{自博弈进化三元角色：Proposer生成任务，Solver尝试求解，Verifier提供反馈信号。三者可由同一模型不同prompt承担。}
\label{fig:ch3-selfplay}
\end{figure}

\subsection{提示词进化}\label{sec:prompt}

提示词进化是最轻量、最容易部署的一类Agent Evolution。\textbf{Reflexion}把稀疏reward转化为"语义梯度"：相比仅告诉模型失败，反思文本能指出失败原因和改进方向。\textbf{ExpeL}从训练任务中自主收集经验，抽取自然语言insights。\textbf{ACE}把上下文工程定义为可演化playbook，使用结构化增量更新避免上下文坍缩。

\subsection{架构搜索}\label{sec:arch}

架构搜索关注agent结构本身的自动设计。\textbf{ADAS}让meta agent用Python编写新agent architecture，在任务上评估后把有效设计加入archive。\textbf{DGM}和\textbf{SICA}把架构搜索推向自我指涉——被改进对象参与改进自身。

\subsection{记忆进化}\label{sec:memory}

记忆进化关注长期经验的积累与检索。\textbf{Voyager}通过自动课程和技能库实现开放世界终身学习。\textbf{Memory-R1}用RL训练记忆管理器（ADD/UPDATE/DELETE/NOOP）。\textbf{AriadneMem}提出演化图表示，用熵感知门控和冲突感知粗化管理记忆。

\subsection{方法对比矩阵}\label{sec:method-compare}

\begin{figure}[htbp]
\centering
\resizebox{\textwidth}{!}{%
\begin{tabular}{l
  >{\centering\arraybackslash}p{22mm}
  >{\centering\arraybackslash}p{22mm}
  >{\centering\arraybackslash}p{16mm}
  >{\centering\arraybackslash}p{16mm}
  >{\centering\arraybackslash}p{16mm}
  >{\centering\arraybackslash}p{16mm}
  >{\centering\arraybackslash}p{20mm}}
\toprule
\textbf{方法类别} & \textbf{进化对象} & \textbf{数据需求} & \textbf{评估可靠性} & \textbf{成本} & \textbf{安全风险} & \textbf{可迁移性} & \textbf{典型系统} \\
\midrule
奖励进化 & 策略/权重 & 中-高 & 中 & 高 & 中 & 中 & STaR, RISE \\
自博弈 & 任务+策略 & 低 & 高 & 中 & 中 & 中 & AbsZero, SPIRAL \\
提示词进化 & 提示/上下文 & 低 & 低-中 & 低 & 低 & 高 & Reflexion, ACE \\
架构搜索 & 代码/结构 & 中 & 中-高 & 高 & 高 & 中 & ADAS, DGM \\
记忆进化 & 经验/技能 & 中 & 中 & 低-中 & 低 & 高 & Voyager, Mem-R1 \\
混合方法 & 多层组合 & 高 & 中-高 & 高 & 中-高 & 中 & EvolveR, WebEvolver \\
\bottomrule
\end{tabular}%
}
\caption{六大方法类别对比矩阵。评估可靠性、成本、安全风险和可迁移性的相对水平。}
\label{fig:ch3-method-compare}
\end{figure}
